\chapter{Complexity — Straight Line Dual Rail}
%%% generated by the blueprint pipeline from TCSlib.Complexity.CircuitComplexity.StraightLineDualRail
\section{Overview}

Declarations of the Lean module \texttt{TCSlib.Complexity.CircuitComplexity.StraightLineDualRail}.

\section{Declarations}

\begin{definition}[Rail]
\lean{BoolCircuit.StraightLine.Rail}
\label{BoolCircuit.StraightLine.Rail}
\leanok
\uses{BoolCircuit.SLOperand}
For $n$ inputs, a rail is an ordered pair of straight-line operands.
\end{definition}

\begin{definition}[Rail at an index]
\lean{BoolCircuit.StraightLine.railAt}
\label{BoolCircuit.StraightLine.railAt}
\leanok
\uses{BoolCircuit.StraightLine.Rail}
Given a list of rails and a natural-number index, return the rail at that index when it
exists; otherwise return the pair consisting of line $0$ twice.
\end{definition}

\begin{definition}[Emit an operation]
\lean{BoolCircuit.StraightLine.emitOp}
\label{BoolCircuit.StraightLine.emitOp}
\leanok
\uses{BoolCircuit.SLOp, BoolCircuit.SLOp.dual, BoolCircuit.SLStmt, BoolCircuit.StraightLine.Rail}
Given an input index $i_0$, an operation $o$ (conjunction or disjunction), a line index
$\ell$, and a list of rails, return a list of statements and an output rail. For an
empty list, emit $o$ applied dually to $x_{i_0}$ and $\neg x_{i_0}$, followed by $o$
applied to those operands, and return the pair of lines $\ell,\ell+1$. For a singleton
list, emit no statements and return its rail. For a list with at least two rails, emit
$o$ on the first components of the first two rails and its dual on their second
components, then return the pair of lines $\ell,\ell+1$.
\end{definition}

\begin{definition}[Emit a gate on rails]
\lean{BoolCircuit.StraightLine.gateEmit}
\label{BoolCircuit.StraightLine.gateEmit}
\leanok
\uses{BoolCircuit.DAGGate, BoolCircuit.SLStmt, BoolCircuit.StraightLine.Rail, BoolCircuit.StraightLine.emitOp}
Given an input index $i_0$, a line index $\ell$, a gate, and a list of rails, return
statements and an output rail as follows. For an AND gate, emit conjunction on the
rails; for an OR gate, emit disjunction. For a NOT gate, emit the same statements as for
conjunction and swap the two components of its output rail.
\end{definition}

\begin{definition}[Dual circuit-processing step]
\lean{BoolCircuit.StraightLine.dualStep}
\label{BoolCircuit.StraightLine.dualStep}
\leanok
\uses{BoolCircuit.DAGGate, BoolCircuit.SLStmt, BoolCircuit.StraightLine.Rail, BoolCircuit.StraightLine.gateEmit, BoolCircuit.StraightLine.railAt}
Given an input index $i_0\in\mathrm{Fin}(n)$, a state consisting of a list of statements
and a list of rails, and a gate $g$, apply the gate-emission procedure to $g$ at the
current statement-list length, using the rails selected by the gate’s argument indices
from the current rail list. Append the emitted statements to the statement list and the
emitted rail to the rail list.
\end{definition}

\begin{definition}[Dual-rail translation]
\lean{BoolCircuit.StraightLine.dualRail}
\label{BoolCircuit.StraightLine.dualRail}
\leanok
\uses{BoolCircuit.DAGGate, BoolCircuit.SLStmt, BoolCircuit.StraightLine.Rail, BoolCircuit.StraightLine.dualStep}
For a natural number $n$, an input index $i_0 in br^n$ (that is, an element of
$	ext{Fin}(n)$), and a list of gates, the dual-rail translation is the pair of statement
and rail lists obtained by folding the gate-by-gate translation over that list. The fold
starts with no statements and, for each input $k$, the rail pair consisting of $x_k$ and
its negation.
\end{definition}

\begin{definition}[Dual-rail circuit invariant]
\lean{BoolCircuit.StraightLine.DualInv}
\label{BoolCircuit.StraightLine.DualInv}
\leanok
\uses{BoolCircuit.DAGGate, BoolCircuit.SLOperand.Before, BoolCircuit.SLOperand.eval, BoolCircuit.SLStmt, BoolCircuit.StraightLine.LinesValid, BoolCircuit.StraightLine.Rail, BoolCircuit.StraightLine.railAt, BoolCircuit.StraightLine.runLines, BoolCircuit.vertexValue}
Let $n$ be a natural number, let $old$ be a finite list of gates, and let $s=(ss,rs)$
consist of a list of straight-line statements and a list of pairs of operands. The rail
list has length $n+|old|$, and each statement reads only input variables or earlier
lines. For every vertex index $v<n+|old|$, both operands in the pair $rs_v$ read only
lines present in $ss$; for every Boolean input assignment, their evaluations are the
value of vertex $v$ in the gate list $old$ and its Boolean complement, respectively. The
statement list has length at most $2|old|$.
\end{definition}

\begin{theorem}[Initial dual-rail invariant]
\lean{BoolCircuit.StraightLine.dualInv_nil}
\label{BoolCircuit.StraightLine.dualInv_nil}
\leanok
\uses{BoolCircuit.DAGGate, BoolCircuit.SLOperand.eval, BoolCircuit.SLStmt, BoolCircuit.StraightLine.DualInv, BoolCircuit.StraightLine.LinesValid.nil, BoolCircuit.StraightLine.Rail, BoolCircuit.StraightLine.railAt, BoolCircuit.vertexValue_input}
\difficulty{4}
For any number $n$ of input variables, start with no gates and no straight-line
statements, and assign to each input vertex $k<n$ the pair consisting of the input $x_k$
and its negation. These rails satisfy the dual-rail invariant: there is one pair per
vertex; all statement and rail references are valid; for every input assignment, the two
rails compute the vertex value and its complement; and the number of statements is at
most twice the number of gates.
\end{theorem}

\begin{theorem}[Evaluation After Appending Two Statements]
\lean{BoolCircuit.StraightLine.runLines_two}
\label{BoolCircuit.StraightLine.runLines_two}
\leanok
\uses{BoolCircuit.SLStmt, BoolCircuit.SLStmt.eval, BoolCircuit.StraightLine.runLines, BoolCircuit.StraightLine.runLines_snoc}
\difficulty{3}
For any natural number $n$, any list $ss$ of straight-line statements on $n$ Boolean
inputs, any two statements $s,t$ on those inputs, and any Boolean input assignment $x$,
evaluating $ss$ followed by $s$ and $t$ gives the evaluation of $ss$ followed by the
values of $s$ on that evaluation and $t$ on the evaluation extended by the value of $s$.
\end{theorem}

\begin{theorem}[Correctness of emitting an operation]
\lean{BoolCircuit.StraightLine.emitOp_spec}
\label{BoolCircuit.StraightLine.emitOp_spec}
\leanok
\uses{BoolCircuit.DAGGate, BoolCircuit.DAGGate.eval, BoolCircuit.SLOp, BoolCircuit.SLOp.apply, BoolCircuit.SLOp.dual, BoolCircuit.SLOp.kind, BoolCircuit.SLOperand.Before, BoolCircuit.SLOperand.Before.mono, BoolCircuit.SLOperand.eval, BoolCircuit.SLOperand.eval_append, BoolCircuit.SLStmt, BoolCircuit.SLStmt.Before, BoolCircuit.SLStmt.eval, BoolCircuit.StraightLine.LinesValid, BoolCircuit.StraightLine.LinesValid.snoc, BoolCircuit.StraightLine.Rail, BoolCircuit.StraightLine.emitOp, BoolCircuit.StraightLine.length_runLines, BoolCircuit.StraightLine.railAt, BoolCircuit.StraightLine.runLines, BoolCircuit.StraightLine.runLines_two}
\difficulty{6}
Let $n$ be the number of input variables, choose an input coordinate and a binary
operation (conjunction or disjunction), and let $ss$ be a valid list of statements with
a list of rails and a list of gate-argument indices. Assume the argument list has at
most two entries, and each referenced rail pair uses only inputs and lines earlier than
$ss$. Emit the operation using the rail pairs at those indices, obtaining new statements
and an output rail pair. The extended statement list is valid, at most two statements
are emitted, and both output rails use only inputs and lines in the extended list.
Moreover, for every input assignment and every list of Boolean vertex values, if each
referenced pair evaluates to the corresponding vertex value and its complement when run
on $ss$, then the output pair, when run on the extended list, evaluates to the indicated
gate's value and its complement; the gate computes the chosen conjunction or disjunction
of the argument values.
\end{theorem}

\begin{theorem}[Correctness of emitting a gate]
\lean{BoolCircuit.StraightLine.gateEmit_spec}
\label{BoolCircuit.StraightLine.gateEmit_spec}
\leanok
\uses{BoolCircuit.DAGGate, BoolCircuit.DAGGate.eval, BoolCircuit.SLOp.kind, BoolCircuit.SLOperand.Before, BoolCircuit.SLOperand.eval, BoolCircuit.SLStmt, BoolCircuit.StraightLine.LinesValid, BoolCircuit.StraightLine.Rail, BoolCircuit.StraightLine.emitOp_spec, BoolCircuit.StraightLine.gateEmit, BoolCircuit.StraightLine.railAt, BoolCircuit.StraightLine.runLines}
\difficulty{4}
Let $n$ be the number of input variables, choose an input coordinate, and let $ss$ be a
valid list of statements with a list of rails and a gate. Suppose the gate has at most
two arguments and each referenced rail pair uses only inputs and lines earlier than
$ss$. Processing the gate using the rail pairs at its argument indices extends the
statements to a valid list, emits at most two statements, and returns a rail pair whose
entries use only inputs and lines in the extended list. For every input assignment and
list of Boolean vertex values, if each argument's rail pair evaluates on the existing
run to that vertex value and its complement, then the output pair on the extended run
evaluates to the gate's value and its complement.
\end{theorem}

\begin{theorem}[Appending Statements to an Evaluation]
\lean{BoolCircuit.StraightLine.runLines_append}
\label{BoolCircuit.StraightLine.runLines_append}
\leanok
\uses{BoolCircuit.SLStmt, BoolCircuit.SLStmt.eval, BoolCircuit.StraightLine.runLines, BoolCircuit.StraightLine.runLines_snoc}
\difficulty{4}
For any natural number $n$, lists $ss$ and $ext$ of straight-line statements on $n$
Boolean inputs, and a Boolean input assignment $x$, there exists a list $e$ such that
evaluating $ss$ followed by $ext$ gives the list obtained by appending $e$ to the
evaluation of $ss$.
\end{theorem}

\begin{theorem}[Operand evaluation after extending a straight-line program]
\lean{BoolCircuit.SLOperand.eval_runLines_append}
\label{BoolCircuit.SLOperand.eval_runLines_append}
\leanok
\uses{BoolCircuit.SLOperand, BoolCircuit.SLOperand.Before, BoolCircuit.SLOperand.eval, BoolCircuit.SLOperand.eval_append, BoolCircuit.SLStmt, BoolCircuit.StraightLine.runLines, BoolCircuit.StraightLine.runLines_append}
\difficulty{3}
Let $x:	ext{Fin}(n)	o	ext{Bool}$ be an assignment, let $ss$ and $ext$ be lists of
straight-line statements, and let $o$ be an operand whose referenced line indices are
all less than the length of $ss$. Then the value of $o$ after running the program formed
by appending $ext$ to $ss$ is the same as its value after running $ss$.
\end{theorem}

\begin{theorem}[Preservation of the dual-rail invariant]
\lean{BoolCircuit.StraightLine.dualInv_step}
\label{BoolCircuit.StraightLine.dualInv_step}
\leanok
\uses{BoolCircuit.DAGGate, BoolCircuit.SLOperand.Before.mono, BoolCircuit.SLOperand.eval_runLines_append, BoolCircuit.SLStmt, BoolCircuit.StraightLine.DualInv, BoolCircuit.StraightLine.Rail, BoolCircuit.StraightLine.dualStep, BoolCircuit.StraightLine.gateEmit, BoolCircuit.StraightLine.gateEmit_spec, BoolCircuit.StraightLine.railAt, BoolCircuit.vertexValue_append, BoolCircuit.vertexValue_last}
\difficulty{6}
Let $n$ be the number of input variables, choose an input coordinate, and suppose a list
of gates and a state of straight-line statements and rails satisfy the dual-rail
invariant. If every argument of a new gate refers to an input or an existing gate
vertex, and the gate has at most two arguments, then after processing that gate the
resulting statements and rails satisfy the dual-rail invariant for the list with the new
gate appended.
\end{theorem}

\begin{theorem}[Correctness of the dual-rail construction]
\lean{BoolCircuit.StraightLine.dualRail_inv}
\label{BoolCircuit.StraightLine.dualRail_inv}
\leanok
\uses{BoolCircuit.DAGGate, BoolCircuit.GatesAcyclic, BoolCircuit.StraightLine.DualInv, BoolCircuit.StraightLine.dualInv_nil, BoolCircuit.StraightLine.dualInv_step, BoolCircuit.StraightLine.dualRail}
\difficulty{5}
Let $n$ be the number of input variables and choose an input coordinate. Suppose a list
of gates is acyclic, meaning that the arguments of its gate at position $i$ all have
index less than $n+i$, and that each gate has at most two arguments. Then the
construction that processes these gates produces statements and rails satisfying the
dual-rail invariant: one rail pair per input or gate vertex, valid references, rails
computing each vertex's value and its complement for every input assignment, and at most
twice as many statements as gates.
\end{theorem}

\begin{definition}[Conversion to a straight-line program]
\lean{BoolCircuit.DAGCircuit.toStraightLine}
\label{BoolCircuit.DAGCircuit.toStraightLine}
\leanok
\uses{BoolCircuit.StraightLine.LinesValid.snoc, BoolCircuit.StraightLine.dualRail, BoolCircuit.StraightLine.dualRail_inv, BoolCircuit.StraightLine.railAt, BoolCircuit.StraightLineProgram}
Let $C$ be a Boolean directed acyclic circuit with $n$ inputs, suppose every gate has at
most two arguments, and choose an input index $i_0\in\mathrm{Fin}(n)$. The associated
straight-line program consists of the statements produced by the dual-rail translation
of the circuit's gate list, followed by one conjunction statement whose two operands are
both the first rail at the circuit's output vertex.
\end{definition}

\begin{theorem}[Evaluation of the straight-line program]
\lean{BoolCircuit.DAGCircuit.toStraightLine_eval}
\label{BoolCircuit.DAGCircuit.toStraightLine_eval}
\leanok
\uses{BoolCircuit.DAGCircuit.eval, BoolCircuit.DAGCircuit.toStraightLine, BoolCircuit.SLOp.apply, BoolCircuit.SLStmt.eval, BoolCircuit.StraightLine.dualRail_inv, BoolCircuit.StraightLine.length_runLines, BoolCircuit.StraightLine.runLines_snoc, BoolCircuit.StraightLineProgram.eval, BoolCircuit.StraightLineProgram.length, BoolCircuit.StraightLineProgram.values}
\difficulty{3}
Let $C$ be a Boolean DAG circuit on $n$ inputs in which every gate has at most two
arguments, and choose an input index $i_0\in\mathrm{Fin}\,n$. For every Boolean input
assignment $x$, the straight-line program constructed from $C$ using $i_0$ has the same
output as $C$.
\end{theorem}

\begin{theorem}[Length bound for the straight-line program]
\lean{BoolCircuit.DAGCircuit.toStraightLine_length_le}
\label{BoolCircuit.DAGCircuit.toStraightLine_length_le}
\leanok
\uses{BoolCircuit.DAGCircuit.toStraightLine, BoolCircuit.StraightLine.dualRail_inv, BoolCircuit.StraightLineProgram.length}
\difficulty{2}
Let $C$ be a Boolean DAG circuit on $n$ inputs in which every gate has at most two
arguments, and choose an input index $i_0\in\mathrm{Fin}\,n$. The length of the
straight-line program constructed from $C$ using $i_0$ is at most
$2|C.\mathrm{gates}|+1$.
\end{theorem}

\begin{theorem}[Circuit representation of a straight-line program]
\lean{BoolCircuit.exists_dagCircuit_of_straightLine}
\label{BoolCircuit.exists_dagCircuit_of_straightLine}
\leanok
\uses{BoolCircuit.DAGCircuit, BoolCircuit.DAGCircuit.IsFaninTwo, BoolCircuit.DAGCircuit.eval, BoolCircuit.DAGCircuit.size, BoolCircuit.StraightLineProgram, BoolCircuit.StraightLineProgram.eval, BoolCircuit.StraightLineProgram.length, BoolCircuit.StraightLineProgram.toCircuit, BoolCircuit.StraightLineProgram.toCircuit_eval, BoolCircuit.StraightLineProgram.toCircuit_isFaninTwo, BoolCircuit.StraightLineProgram.toCircuit_size}
\difficulty{3}
Let $n,S$ be natural numbers, and suppose a straight-line program on $n$ Boolean inputs
has length at most $S$ and computes a Boolean function $f$. Then there is a well-formed
acyclic Boolean circuit with at most two inputs to each gate that computes $f$ and has
size at most $2n+\max(S,1)$.
\end{theorem}

\begin{theorem}[Straight-line simulation of a fan-in-two circuit]
\lean{BoolCircuit.exists_straightLine_of_dagCircuit}
\label{BoolCircuit.exists_straightLine_of_dagCircuit}
\leanok
\uses{BoolCircuit.DAGCircuit, BoolCircuit.DAGCircuit.IsFaninTwo, BoolCircuit.DAGCircuit.eval, BoolCircuit.DAGCircuit.size, BoolCircuit.DAGCircuit.toStraightLine, BoolCircuit.DAGCircuit.toStraightLine_eval, BoolCircuit.DAGCircuit.toStraightLine_length_le, BoolCircuit.StraightLineProgram, BoolCircuit.StraightLineProgram.eval, BoolCircuit.StraightLineProgram.length}
\difficulty{3}
Let $n>0$, and suppose a Boolean function on $n$ inputs is computed by a fan-in-two
directed acyclic circuit of size at most $S$. Then the function is computed by a valid
straight-line program of length at most $2(S-n)+1$.
\end{theorem}

\begin{theorem}[A constant true function separates circuit models]
\lean{BoolCircuit.not_exists_straightLine_true}
\label{BoolCircuit.not_exists_straightLine_true}
\leanok
\uses{BoolCircuit.DAGCircuit, BoolCircuit.DAGCircuit.IsFaninTwo, BoolCircuit.DAGCircuit.eval, BoolCircuit.DAGCircuit.size, BoolCircuit.StraightLineProgram, BoolCircuit.StraightLineProgram.eval, BoolCircuit.StraightLineProgram.eval_of_zero, BoolCircuit.constCircuit, BoolCircuit.constCircuit_eval, BoolCircuit.constCircuit_isFaninTwo}
\difficulty{4}
There exists a well-formed acyclic Boolean circuit with fan-in at most two, no input
variables, and size one that evaluates to true for every input assignment. There is no
valid straight-line program with no input variables that evaluates to true for every
input assignment.
\end{theorem}
